\documentclass[12pt,a4paper]{article} % using article ensures it starts at 1 and does not have odd numberings for section
\usepackage{graphicx}
\usepackage[utf8]{inputenc} % this way umlaute are included from the get go
\usepackage[ngerman]{babel} % german spell check
\usepackage{datetime}

\usepackage{breqn} % this package is one option for math lines

\usepackage{hyperref} % these two lines are so that the table of content is clickable
\hypersetup{linktoc=all}

\usepackage{amssymb} % package for Natural Number sign etc

\newcommand{\itab}[1]{\hspace{0em}\rlap{#1}}
\newcommand{\tab}[1]{\hspace{.2\textwidth}\rlap{#1}} % own tab function

\newdateformat{gerDate}{\THEDAY \space \monthname[\THEMONTH], \THEYEAR} % german date format

\begin{document}
	\begin{titlepage} % good tital page template, only needs the class X notes part to be changed for each new class
	\centering
	\includegraphics[width=0.40\textwidth]{UniLogo}\par\vspace{1cm}
		{\scshape\LARGE Universität Heidelberg \par}
		\vspace{1cm}
		{\huge\bfseries Numerik Zusammenfassung \par}
		\vspace{2cm}
		{\Large\itshape by Charles Barbret \par}
		
		\vfill

		% Bottom of the page
		{\large \gerDate\today\par}
	\end{titlepage}

	\tableofcontents % creats a table of contents, ensured already that it is clickable
	\newpage % starts the actual document on a new page so there is no weird colision of text and toc
	
	
	\section{Fehleranalyse}
	\subsection{}
	Zahlendarstellung und Rundungsfehler Schema für Gleitkomma Zahl:
	\begin{dmath}
		x = \pm m * b^{\pm e}
	\end{dmath}
	Basis $b \in \mathbb{N}$ \space b $\geq 2 \space $ \footnote{Beispiel: $ 2^e$ oder $10^e $}
	\newline
	Mantisse $ m = m_1 b^{-1} + m_2 b^{-2} + ... \in \mathbb{R} $ \footnote{ Beispiel: $ m_1 = 3, m_2 = 1, m_3 = 4 \Rightarrow m = 314 $ normal kommt nach $ m_1 $ ein Komma, also m = 3,14}
	\newline
	Exponent e = $e_1 b^{1} + e_2 b^{2} + ... \in \mathbb{N}_0 $
	\newline
	$ \forall m_i \in m $ und $ \forall e_i \in e $ gilt $e_i , m_i \in $ \{0, ..., b - 1\}
	\newline
	Sollte b = 10 sein befinden wir uns im Dezimalsystem
	\newline \indent 
	$\Rightarrow$ Es gibt keine Ziffer > 9 = 10 - 1 $\Leftrightarrow$ b - 1 $q.e.d$
	\newline
	Jede Ziffer ($m_i , e_i , b, \pm, \pm$) braucht man eine Speicherzelle, wobei b im Computer bereits eingespeichert ist
	\newline \indent 
	$\Rightarrow $ muss nicht explizit angegeben werden
	\newline
	X wird gespeichert als: ($\pm$)[$m_1 , ..., m_r $]($\pm$)[$e_{s - 1} , ..., e_0 $]
	
\end{document}